\begin{python0}
from solutions import *; clear()
\end{python0}
\sectionthree{Summary}

The if and if-else statement looks like these:
\begin{align*}
&\text{if \texttt{([bool expr])}}\\
&\tab[3em]{\texttt{[stmt]}}\\\\
&\text{if \texttt{([bool expr])}}\\
&\tab[3em]{\texttt{[stmt1]}}\\
&\texttt{else}\\
&\tab[3em]{\texttt{[stmt2]}}\\
\end{align*}

where \texttt{[bool expr]} is a boolean expression and
\texttt{[stmt], [stmt1]}, and \texttt{[stmt2]} are either
statements or blocks of statements.

The else statement will always be associated with the closest if
statement. In other words
\begin{console}
if (x < 0)
   if (x == -2) std::cout << "1" << std::endl;
else
      std::cout << "2" << std::endl;
\end{console}

is better written this way:
\begin{console}
if (x < 0)
   if (x == -2)
       std::cout << "1" << std::endl;
   else
       std::cout << "2" << std::endl;
\end{console}

The \texttt{rand()} function returns a \lq\lq random'' integer between 0 and
\texttt{RAND\_MAX}. Number. The \texttt{srand()} function \lq\lq seeds'' the
\texttt{rand()} function. You might need to \lq\lq \texttt{\#<include
 cstdlib>}'' to use the \texttt{srand()} and
\texttt{rand()} functions. One way to seed rand() is to call
\begin{center}
\texttt{srand((unsigned) time(0))}
\end{center}
in which case you need to \lq\lq \texttt{\#<include  ctime>}''.
\begin{center}
\texttt{double(rand()) / RAND\_MAX}
\end{center}
will generate \lq\lq random'' doubles between 0 and 1. To generate a random
integer between a and b (two integers with a = b), you can use
\begin{center}
\texttt{int(double(rand()) / RAND\_MAX * (b - a)) + a}
\end{center}
or
\begin{center}
\texttt{rand() \% (b - a + 1) + a}
\end{center}
(The second method is not as general nor as random as the first. For
instance you cannot use it to generate random number between 0 and
1000000 in MS .NET STUDIO).

The following code segment swaps the values in \texttt{a} and \texttt{b}
(and \texttt{t} has the value of \texttt{a}):
\begin{align*}
&\text{t = a;} \tab[3em]{\text{a = b;}} \tab[3em]{\text{b = t;}}\\
\end{align*}
Using the \lq\lq swapping trick'', you can always ensure that a pair of
variables are in ascending (or descending order):

\tab[3em]{if (a $>$ b)}\\
\tab[3em]{\{}\\
\tab[6em]{t = a;}\\
\tab[6em]{a = b;}\\
\tab[6em]{b = t;}\\
\tab[3em]{\}}\\

Given for instance four variables \texttt{a}, \texttt{b}, \texttt{c}, and
\texttt{d}, you can apply the above to the pairs \texttt{a},\texttt{b} and
\texttt{b},\texttt{c} and \texttt{c},\texttt{d} to ensure that the largest value
among \texttt{a}, \texttt{b}, \texttt{c}, \texttt{d} is in \texttt{d:}

\tab[3em]{if (a $>$ b) \{ swap values of a and b \}}\\
\tab[3em]{if (b $>$ c) \{ swap values of b and c \}}\\
\tab[3em]{if (c $>$ d) \{ swap values of c and d \}}\\

This is called one pass of the sorting process. You can then apply the
process to a, b, c for a second pass, and then to a, b for a third pass.
This will sort the values of \texttt{a}, \texttt{b}, \texttt{c}, \texttt{d} into
ascending order.

Just as there is automatic type casting from int to double, there is
also automatic type type casting between bool and int or double:

\begin{align*}
&\texttt{bool} \implies &\texttt{int}\\
&\texttt{true} &1\\
&\texttt{false} &0\\\\
&\text{numeric (int/double/float)} \implies &\texttt{bool}\\
&\text{nonzero (example: 1, 2, 3)} &\texttt{true}\\
&\text{zero (example: 0, 0.0)} &\texttt{false}\\
\end{align*}

For instance the following is a valid C++ statement:
\tab[3em]{if (1)}
\tab[6em]{std::cout $<<$ "that's all folks";}

The ternary operator looks like this:
\begin{center}
(\texttt{[bool expr]} ? \texttt{[expr1]} : \texttt{[expr2]})
\end{center}
This whole expression evaluates to \texttt{[expr1]} if \texttt{[bool
expr]} is true; otherwise it evaluates to \texttt{[expr2]}.

